\section{Detailed derivations and limiting cases}
\label{app:derivations}

\subsection{From quantum relative entropy to a product free energy}

The Gibbs target of Eq.~\eqref{eq:gibbs} is full rank for $T>0$. Starting
from the definition of relative entropy,
\begin{align}
 \KL(\rho\Vert\sigma_{T,\Gamma})
 &=\Tr(\rho\log\rho)-\Tr(\rho\log\sigma_{T,\Gamma})\nonumber\\
 &=-S(\rho)+\frac{\Tr(\rho H_\Gamma)}{T}
                 +\log Z_{T,\Gamma}.
 \label{eq:app-kl}
\end{align}
Multiplication by $T$ immediately gives Eq.~\eqref{eq:klidentity}.
The omitted constant in an optimization is $-T\log Z$, which depends
on $(T,\Gamma)$ but not on the variational fields. Omitting it from
gradient calculations is valid even though its value changes along the
schedule. The exact Gibbs minimum is outside the factorized family in
general, so the restricted minimum only bounds the exact free energy
from above.

For a product state $\rho=\bigotimes_i\rho_i$, the logarithm separates
as $\log\rho=\sum_i I\otimes\cdots\otimes\log\rho_i\otimes\cdots\otimes I$,
giving $S(\rho)=\sum_iS(\rho_i)$. For $i\ne j$,
\begin{equation}
 \Tr[\rho(Z_iZ_j)]
 =\Tr(\rho_iZ_i)\Tr(\rho_jZ_j)=m_{z,i}m_{z,j}.
 \label{eq:app-factor}
\end{equation}
The transverse expectation is $\Tr(\rho X_i)=m_{x,i}$. These identities
give Eq.~\eqref{eq:qefemF} without sampling any spin configuration.

\subsection{The local state, entropy, and field map}

Let $A_i=u_iX_i+v_iZ_i$. The Pauli anticommutation relation gives
$A_i^2=(u_i^2+v_i^2)I=r_i^2I$. Separating even and odd terms of the
exponential series,
\begin{equation}
 e^{A_i}=\cosh r_i\,I+
           \frac{\sinh r_i}{r_i}A_i.
 \label{eq:app-exp}
\end{equation}
Dividing by $\Tr e^{A_i}=2\cosh r_i$ produces
$\rho_i=\tfrac12[I+(\tanh r_i/r_i)A_i]$. Thus the Bloch radius is
$R_i=\tanh r_i$, and the eigenvalues are $(1\pm R_i)/2$.
Substituting those eigenvalues into $-\Tr\rho_i\log\rho_i$ yields
Eq.~\eqref{eq:entropy}. Alternatively, from
$\log\rho_i=A_i-\log(2\cosh r_i)I$,
\begin{align}
 S_i&=-\Tr(\rho_i A_i)+\log(2\cosh r_i)\nonumber\\
    &=\log(2\cosh r_i)-r_i\tanh r_i,
 \label{eq:app-entropy}
\end{align}
because $\Tr(\rho_i A_i)=r_iR_i$.

The apparent singularity at $r=0$ in $g(r)=\tanh r/r$ is removable.
Expanding $\tanh r$ around zero gives
\begin{equation}
 g(r)=1-\frac{r^2}{3}+\frac{2r^4}{15}+O(r^6),
 \qquad
 \operatorname{sech}^2r=1-r^2+\frac{2r^4}{3}+O(r^6).
 \label{eq:app-series}
\end{equation}
These limits justify a smooth origin and the numerical safe branch in
the implementation.

\subsection{Gradient and Hessian in physical and field coordinates}

Writing $\bm m=(m_x,m_z)$ and $R=\|\bm m\|$, the scalar derivative of
Eq.~\eqref{eq:entropy} is
\begin{equation}
 \frac{dS}{dR}
 =-\frac12\log\frac{1+R}{1-R}
 =-\operatorname{arctanh}R.
 \label{eq:app-ds}
\end{equation}
As $\nabla_{\bm m}R=\bm m/R$ away from the origin,
$\nabla_{\bm m}(-TS)=T\operatorname{arctanh}R\,\bm m/R$.
Adding the problem and driver gradients gives
Eqs.~\eqref{eq:gradx}--\eqref{eq:gradz}. To move to unconstrained
fields, differentiate $\bm m=g(r)\bm a$:
\begin{equation}
 d\bm m=g(r)d\bm a+
   g'(r)\bm a\frac{\bm a^{\mathsf T}d\bm a}{r}.
 \label{eq:app-dm}
\end{equation}
This is Eq.~\eqref{eq:jacobian}. In the radial direction
$\bm a/r$, the Jacobian eigenvalue is
$g(r)+rg'(r)=d(\tanh r)/dr=\operatorname{sech}^2r$;
in every tangential direction it is $g(r)$. Since
$J\bm a=\operatorname{sech}^2r\,\bm a$,
the chain rule applied to Eq.~\eqref{eq:physicalgrad} gives
Eq.~\eqref{eq:fieldgrad}. No derivative of a sign function appears in
this exact smooth calculation.

For completeness, the physical-coordinate Hessian of the negative
single-site entropy has the radial and tangential eigenvalues
\begin{equation}
 \lambda_{\rm rad}=\frac{1}{1-R^2},
 \qquad
 \lambda_{\rm tan}=\frac{\operatorname{arctanh}R}{R}.
 \label{eq:app-entropy-hessian}
\end{equation}
They follow by differentiating the radial vector field
$\nabla(-S)=(\operatorname{arctanh}R/R)\bm m$. Since
$\operatorname{arctanh}R\geq R$ for $0\leq R<1$, both are at least one.
The graph contribution to the Hessian acts only on longitudinal
coordinates and equals $K$. Thus $T> -\lambda_{\min}(K)$ is a sufficient
condition for strict convexity of the physical-coordinate objective.
This does not imply convexity in the nonlinear field coordinates
$\bm a$; reparameterization may change curvature even while preserving
finite-field stationary points.

\subsection{Longitudinal instability and spectral schedule inversion}

When $h=0$ and $\bm m_z=0$, the stationary local $x$ component is
$m_x=\tanh(\Gamma/T)$. The graph term has quadratic expansion
$\tfrac12\bm m_z^{\mathsf T}K\bm m_z$. Holding the stationary transverse
component fixed for a longitudinal perturbation, the local entropy
curvature contributes
\begin{equation}
 T\frac{\operatorname{arctanh}m_x}{m_x}
 =\frac{\Gamma}{\tanh(\Gamma/T)}=q(T,\Gamma).
 \label{eq:app-q}
\end{equation}
The mixed $x$--$z$ Hessian vanishes at $m_z=0$ by symmetry, so
Eq.~\eqref{eq:stability} is also the relevant longitudinal block of
the full Hessian. If the lowest eigenvalue of $K$ is $-\kappa$, the
quadratic form first loses positive definiteness at $q=\kappa$.
Setting $\rho=\Gamma/T$ gives
$q=T\rho/\tanh\rho$. Algebraic inversion yields
$T=q\tanh\rho/\rho$ and $\Gamma=q\tanh\rho$, as used in the
graph-specific G-set schedule. The stability calculation assumes zero
local fields and concerns a local stationary solution; it is not a
global approximation guarantee.

\subsection{Quadratization example and its penalty requirement}

The formal reach of a QUBO solver extends to higher-order binary terms
when auxiliary variables are introduced, although that extension is not
benchmarked here. For $x,y,z\in\{0,1\}$, the penalty
\begin{equation}
 P\bigl(xy-2xz-2yz+3z\bigr),\qquad P>0,
 \label{eq:app-quadratize}
\end{equation}
is zero when $z=xy$ and strictly positive for the other assignments.
For example, $(x,y)=(0,0)$ gives penalty $3Pz$;
$(1,0)$ or $(0,1)$ gives $Pz$; and $(1,1)$ gives $P(1-z)$.
To preserve the optimum of a surrounding objective, $P$ must exceed
the maximum possible improvement gained by violating the intended
relation. Repeated substitutions can increase variable count and
coefficient spread. Consequently QUBO encodability alone is not a
performance claim for a new application.
